\subsection{Two auxiliary units from a Jacobi heat lift}
\label{sec:kt-jacobi-fields}

The following construction varies the imaginary quadratic field.  Its
extra rational Bernoulli condition is $B_{2b}/(2b)$, whereas the
original rational reflection component at $e\equiv b-1$ is controlled
by $B_{p-b}/(p-b)$.  No unit condition on the latter number is assumed.
The construction does not assert simultaneous units for both twists
$\omega^e$ and $\omega^{2e}$ over the same field.

\begin{lemma}[KT-SMALL-DISCRIMINANT-CLASS-BOUND]
Let $D<0$ be an odd fundamental discriminant, put $d=|D|$ and
$R=\lfloor\sqrt{d/3}\rfloor$.  Then
\[
 h(D)\leq\frac{R(R+3)}2
       \leq\frac d6+\frac{\sqrt{3d}}2.                 \tag{JF.1}
\]
Consequently, if $p\geq29$ and $d<4p$, then $h(D)<p$.
\end{lemma}
\begin{proof}
Here is the elementary form-counting argument, including the
representation of ideal classes that the count uses.  For an integral
ideal $I$ and an oriented integral basis $\alpha,\beta$, the form
\[
 Q_I(x,y)=N(x\alpha+y\beta)/N(I)=ax^2+bxy+cy^2
\]
is integral and positive definite, and has discriminant $D$.
Integrality follows because $(u)\subseteq I$ implies $N(I)\mid N(u)$,
also applied to $u=\alpha+\beta$.  The discriminant follows from
the covolume $N(I)\sqrt d/2$.  An oriented basis change is a proper
equivalence.  Conversely, a proper equivalence between two such norm
forms gives an orientation-preserving similarity between their complex
ideal lattices, hence multiplication by some $\lambda\in\mathbf C$.
Since it maps a nonzero element of one ideal to an element of the other,
$\lambda\in\mathbf Q(\sqrt D)$; the two ideals have the same class.
Thus ideal classes inject into proper equivalence classes of these forms.

In each form class choose a primitive vector with smallest positive
value $a$ and extend it to an oriented integral basis.  The vector is
primitive because otherwise division by its common divisor would
decrease its positive value.  Changing the second basis vector by an
integer multiple of the first arranges $|b|\leq a$.  Its value $c$ is
at least $a$ by minimality.  Hence $d=4ac-b^2\geq3a^2$ and $a\leq R$.
For a fixed $a$, the odd integer $b$ lies in $[-a,a]$, giving at most
$a+1$ choices; $c=(b^2+d)/(4a)$ is then determined.  Counting all these
pairs, even the ones that are not integral forms, proves (JF.1).
If $d<4p$, its right side is less than $2p/3+\sqrt{3p}$, which is
less than $p$ when $p>27$.
\end{proof}

\begin{lemma}[KT-JACOBI-ODD-COEFFICIENT]
Let $p\geq5$, let $K$ be even, and let
$\Phi=\sum a(n,r)q^n\zeta^r\in J_{K,1}(\mathrm{SL}_2(\mathbf Z))$
have rational $p$-integral coefficients with bounded denominators.
If $\Phi\not\equiv0\pmod p$, then its odd-$r$ theta component is
nonzero modulo $p$.  Moreover, some coefficient $a(n,r)$ with $r$ odd
and $0<n\leq K/2$ is a $p$-adic unit.
\end{lemma}
\begin{proof}
Use $e(x)=\exp(2\pi i x)$ and the usual Jacobi slash operator.
Put $g(\tau,z)=\Phi(\tau,z+1/2)$ and $F=\Phi-g$.
The elliptic transformation law shows that $g$ has index one for
the full integral translation lattice.  For
$\gamma=\left(\begin{smallmatrix}a&b\\c&d\end{smallmatrix}\right)
\in\Gamma_0(4)$, the modular law followed by the integral translation
$(c/2,(d-1)/2)$ gives exactly
\[
                   g|_{K,1}\gamma=e(c(d-2)/4)g=g.       \tag{JF.2}
\]
Thus $F$ is a Jacobi form for $\Gamma_0(4)$, and in particular for
$\Gamma_1(2^2)\ltimes(2\mathbf Z\times\mathbf Z)$.
Holomorphy at its cusps follows directly by transforming to rational
elliptic translates: the resulting exponents are
$n+r\lambda+\lambda^2=(n-r^2/4)+(\lambda+r/2)^2\geq0$.

Suppose all odd-$r$ coefficients of $\Phi$ vanish modulo $p$.
Then $F\equiv0\pmod p$.  Schwagenscheidt, Theorem~1.2
(statement p.~2, complete proof pp.~11--12), applied with $N=2,m=1$,
implies $F|_{K,1}S\equiv0\pmod p$ for
$S=\left(\begin{smallmatrix}0&-1\\1&0\end{smallmatrix}\right)$.
Its integrality hypothesis is met by multiplying by a common
denominator prime to $p$.  An explicit calculation gives
\[
 F|_{K,1}S=\Phi(\tau,z)
       -q^{1/4}\zeta\Phi(\tau,z+\tau/2).                \tag{JF.3}
\]
Modulo $p$, the first summand has only even powers of $\zeta$ and
the second only odd powers.  Their vanishing difference forces
$\Phi\equiv0$, a contradiction.  This proves the first assertion.

Write the index-one theta decomposition as
\[
 \Phi=h_0\vartheta_0+h_1\vartheta_1,\qquad
 \vartheta_\mu(\tau,z)=
       \sum_{r\equiv\mu\ (2)}q^{r^2/4}\zeta^r.
\]
Then $h_1\not\equiv0\pmod p$, and
\[
 f(\tau)=\frac{\Phi(\tau,0)-\Phi(\tau,1/2)}2
        =h_1(\tau)\vartheta_1(\tau,0)
        \in M_K(\Gamma_0(4)).                           \tag{JF.4}
\]
The character is trivial, by (JF.2) at $z=0$; there is no unrecorded
multiplier.  The series $\vartheta_1(\tau,0)$ starts with $2q^{1/4}$,
so (JF.4) is nonzero modulo $p$; cancellation between the odd-$r$
coefficients cannot make it zero.  It is $p$-integral.
The ordinary Sturm bound gives a nonzero coefficient of $f$ at
$n\leq K[\mathrm{SL}_2(\mathbf Z):\Gamma_0(4)]/12=K/2$.
This is the usual elliptic Sturm theorem, stated in exactly this
coefficient-field generality on p.~129 of Choi--Choie--Kikuta (2013),
with the original reference Sturm (1987), pp.~275--280.
That coefficient is $\sum_{r\text{ odd}}a(n,r)$, so one of its
summands is a unit.  The Jacobi support condition rules out $n=0$.
\end{proof}

\begin{theorem}[KT-JACOBI-TWO-UNIT-FIELD]
Let $p\geq29$ be prime and let $b$ be odd with
\[
 3\leq b\leq p-2,\qquad b\ne(p-1)/2,
 \qquad B_{2b}/(2b)\in\mathbf Z_p^\times.               \tag{JF.5}
\]
There exists an odd negative fundamental discriminant $D$, prime to
$p$, such that
\[
 |D|\leq2(b+p+2)-1\leq4p-1,\qquad
 B_{1,\chi_D},\quad B_{1,\chi_D\omega^{b-1}}
                      \in\mathbf Z_p^\times.           \tag{JF.6}
\]
Equivalently, $p\nmid h(D)$ and
$B_{b,\chi_D}/b\in\mathbf Z_p^\times$.
The same field has $B_{n,\chi_D}/n$ a unit for every odd $n\geq3$
congruent to $b$ modulo $p-1$, with no bound on $n$.
\end{theorem}
\begin{proof}
Put $k=b+1$ and take the normalized Jacobi Eisenstein series
\[
                \phi=\zeta(1-2b)E_{b+1,1}\in J_{k,1}.
                                                               \tag{JF.7}
\]
The index-one coefficient formula of Eichler--Zagier, Theorem~2.1,
is also recalled for fundamental discriminants by Schwagenscheidt
on p.~3.  Its general-discriminant version is explicit in the
1983 MPIM preprint: the definition of $L_D(s)$ on printed p.~19a
(PDF p.~24), followed by Theorem~2.1 on printed p.~20 (PDF p.~25).
In full, if $r^2-4n=D_0u^2<0$ with $D_0$ fundamental,
then the coefficient of (JF.7) is
\[
 L(1-b,\chi_{D_0})
 \sum_{v\mid u}\mu(v)\chi_{D_0}(v)v^{b-1}
                      \sigma_{2b-1}(u/v).               \tag{JF.8}
\]
For zero discriminant it is $\zeta(1-2b)$.
The coefficient denominators are bounded; this property of (JF.7)
is also stated on p.~3 of Schwagenscheidt.

We verify $p$-integrality even when $p\mid D_0$, so this step does
not depend on an omitted condition on the discriminant.  For any
quadratic primitive character $\chi$ of conductor $f$, the Bernoulli
polynomial formula is
\[
 \frac{B_{b,\chi}}b
 =\frac1b\sum_{j=0}^{b}\binom bj B_j f^{j-1}
                     \sum_{a=1}^{f}\chi(a)a^{b-j}.       \tag{JF.9}
\]
Since $b\leq p-2$, every $B_j$ here is $p$-integral, as follows
from its elementary recurrence, and $b$ is a unit.  If $p\nmid f$
the assertion follows immediately.  If $p\mid f$, then $v_p(f)=1$
because $f$ is a fundamental discriminant and $p$ is odd.
Only the $j=0$ term can have a denominator $p$.  Decompose the
primitive quadratic character by the Chinese remainder theorem
as $\chi_p\chi_{f/p}$.  Modulo $p$, its numerator is a product
of the sums of $\chi_p(a)a^b$ over $a\in\mathbf F_p^\times$
and of $\chi_{f/p}$ over its complete period.  The latter sum is
zero unless $f/p=1$.  In that case the first sum is
$\sum_{a\in\mathbf F_p^\times}a^{b+(p-1)/2}$, which is zero
unless $b=(p-1)/2$.  That exception was excluded in (JF.5).
Thus (JF.9), and consequently every coefficient in (JF.8), is
$p$-integral.  The zero-discriminant coefficient is a unit by (JF.5).
In particular $\phi\not\equiv0\pmod p$.

Richter (2009), Proposition~1 on p.~872, gives the weight
congruence for nonzero Jacobi forms modulo $p$.
It implies that the filtration of $\phi$ is $k$, since
$4\leq k\leq p-1$.  At $k=p-1$ the only lower possible weight
is zero, and $J_{0,1}=0$; this also follows from the free
$E_{4,1},E_{6,1}$ decomposition recalled on p.~871 of Richter.
Proposition~2, including its proof on p.~873, now gives
\[
 \Omega(L_1\phi)=k+p+1=:K,\qquad L_1\phi\not\equiv0\pmod p,
                                                               \tag{JF.10}
\]
because $p\nmid 2k-1=2b+1$.  In the indicated range the only
failure of this last condition is the same excluded $b=(p-1)/2$.
For precision an actual lift is
\[
 \widehat\phi=L_1\phi-\frac{2k-1}{6}E_2\phi\in J_{k+2,1},
 \qquad
 \Phi=E_{p-1}\widehat\phi+
                   \frac{2k-1}{6}E_{p+1}\phi\in J_{K,1}.
                                                               \tag{JF.11}
\]
Here $\Phi\equiv L_1\phi\pmod p$, using
$E_{p-1}\equiv1$ and $E_{p+1}\equiv E_2$; all coefficients
are $p$-integral.  Formula (JF.11) is the explicit lift used in
Richter's proof, rather than merely a formal heat series.

Apply KT-JACOBI-ODD-COEFFICIENT to $\Phi$.  It gives a unit
coefficient at $(n,r)$ with $r$ odd and $n\leq K/2$.
Modulo $p$ that coefficient equals
$(4n-r^2)c_\phi(n,r)$.  Therefore $p\nmid4n-r^2$, and
$c_\phi(n,r)$ is itself a unit.  The integer
$D'=r^2-4n$ is negative and odd, and
\[
                    0<|D'|\leq4n-1\leq2K-1\leq4p-1.
                                                               \tag{JF.12}
\]
Write $D'=Du^2$ with $D$ fundamental.  Then $D$ is odd,
$p\nmid Du$, and $|D|\leq|D'|$.
The multiplier in (JF.8) is an integer.  Since
$L(1-b,\chi_D)=-B_{b,\chi_D}/b$ is $p$-integral,
the unit coefficient forces $B_{b,\chi_D}/b$ to be a unit.
KT-SMALL-DISCRIMINANT-CLASS-BOUND gives $h(D)<p$, hence the
first unit in (JF.6), by $B_{1,\chi_D}=-2h(D)/w(D)$.
Finally KT-FIELD-CHARACTER-REDUCTION, (FV.2), gives the second
unit in (JF.6) and the asserted transport to all weights $n$.
\end{proof}

\begin{corollary}[KT-JACOBI-TWO-UNIT-GD]
Under (JF.5), suppose also $(p-1)\nmid2(b-1)$ and let
$e=b-1+(p-1)t$, $t\geq0$.  The field supplied by (JF.6) gives
actual \texttt{GD} data for $N=2$, $\psi=1$, and this $e$.
It has the class-number unit and the unit
$B_{1,\chi_D\omega^e}$.

In particular, for every prime $p\geq31$ with $p\equiv3\pmod4$
and every
\[
                      e=(p+1)/2+(p-1)t,\qquad t\geq0,    \tag{JF.13}
\]
such an actual GD field exists without any additional rational
Bernoulli hypothesis.  Here $\omega^{1+e}$ has order $p-1$,
and $\omega^e$ has order $(p-1)/2$.
\end{corollary}
\begin{proof}
The twelve clauses of \texttt{GD} are the ones in
\texttt{Counts.lean}, lines 9--22.  The verification in
KT-ODD-DISCRIMINANT-GD applies: the discriminant is odd and
prime to $p$, its class number is prime to $p$, $e$ is positive
and even, and the required nondivisibility is explicitly assumed.
For $\psi=1$ the order and square-compatibility clauses are trivial;
$N=2$ is even and coprime to $D$.  The quadratic and imaginary
clauses hold by construction.  For the remaining \texttt{eta}
clause, $E$ cannot be contained in $\mathbf Q(\mu_p)$: some
prime dividing $D$ ramifies in $E$ and is different from $p$,
whereas $\mathbf Q(\mu_p)$ is unramified there.  This also makes
explicit that no splitting-at-$p$ assumption is used.

For (JF.13), set $b=(p+3)/2$.  It is odd, is in the required
range, and is not $(p-1)/2$.  The ordinary Kummer congruence
(the already checked Sun lemma with trivial character) gives
\[
 \frac{B_{2b}}{2b}=\frac{B_{p+3}}{p+3}
       \equiv\frac{B_4}{4}=-\frac1{120}\pmod p.          \tag{JF.14}
\]
This is a unit.  Also $2e\equiv2\pmod{p-1}$.
Writing $q=(p-1)/2$, which is odd, gives
$\gcd(2q,q+2)=1$ and $\gcd(2q,q+1)=2$.
These are exactly the stated character orders.  In particular
the mirror in (JF.13) is nonquadratic; no vanishing of its
rational class component was an input to the construction.
\end{proof}

\begin{corollary}[KT-JACOBI-TYPEII-FIELD]
Let $p\geq31$, let $e\geq2$ be even with $(p-1)\nmid2e$,
and let $b$ be the odd representative in $[3,p-2]$ of $2e+1$.
If $b\ne(p-1)/2$ and $B_{2b}/(2b)$ is a $p$-adic unit,
there is an odd imaginary quadratic GD field for $N=2$, $\psi=1$
over which the actual source has no TypeII points.
In particular this holds, without an extra Bernoulli hypothesis, for
\[
 p\equiv7\pmod8,\qquad
                e=(p+1)/4+(p-1)t,\qquad t\geq0.        \tag{JF.15}
\]
\end{corollary}
\begin{proof}
Apply KT-JACOBI-TWO-UNIT-FIELD to this $b$.  Its field has
$p\nmid h(D)$ and $B_{1,\chi_D\omega^{2e}}$ a unit.
The same twelve-clause verification just given supplies GD for
the specified $e$, whose genericity was assumed.  The previously
proved KT-AUXILIARY-UNIT-NO-TYPEII applies to these actual source
data, including its finite-coefficient descent and full rational
ordinary flag; it gives the conclusion about TypeII points.
For (JF.15), $e$ is even, $2e+1\equiv(p+3)/2$, and (JF.14)
supplies the seed.  Here $2e\equiv(p+1)/2\not\equiv0\pmod{p-1}$.
This corollary does not assert that $B_{1,\chi_D\omega^e}$ is a
unit for the same field, and does not exclude genuine rank-three
sources.
\end{proof}

\paragraph{Primary sources and exact scope.}
M.~Eichler and D.~Zagier, \emph{On the theory of Jacobi forms I},
MPIM Preprint 1983(8), printed pp.~19a--20 (PDF pp.~24--25),
gives the general-discriminant $L_D$ definition and Theorem~2.1
used in (JF.8).  The public primary scan is
\url{https://archive.mpim-bonn.mpg.de/id/eprint/1437/1/preprint_1983_8.pdf}.
It is a precursor of their 1985 monograph \emph{The Theory of Jacobi
Forms}, Progress in Mathematics~55; the displayed page locators here
refer to the 1983 preprint, not to the monograph.
O.~K.~Richter, \emph{The action of the heat operator on Jacobi forms},
Proc. Amer. Math. Soc.~137 (2009), 869--875,
doi:10.1090/S0002-9939-08-09566-X, provides the free decomposition
on p.~871, Proposition~1 on p.~872, and Proposition~2 with its
complete proof on p.~873.
M.~Schwagenscheidt, \emph{Nonvanishing modulo $\ell$ of Fourier
coefficients of Jacobi forms}, Research in Number Theory~2 (2016),
article~4, doi:10.1007/s40993-015-0035-1, supplies Theorem~1.2
(p.~2, full proof pp.~11--12) and the normalized Eisenstein and
bounded-denominator statements on p.~3.
D.~Choi, Y.~Choie and T.~Kikuta, \emph{Sturm type theorem for
Siegel modular forms of genus~2 modulo $p$}, Acta Arith.~158
(2013), 129--139, doi:10.4064/aa158-2-2, states the ordinary
elliptic Sturm theorem on p.~129, citing J.~Sturm,
\emph{On the congruence of modular forms}, Lecture Notes in
Mathematics~1240 (1987), 275--280.  The ordinary Sturm theorem is
imported through that exact published statement; no fresh reading
of Sturm's original proof is claimed.
